\begin{scanpage}{0240}{222}
{\cyrillicfont\small
\noindent сфере. «Радиотехника и электроника», 3, № 12, 1451--1462, 1958.\par
75. Chytil B. The depolarisation of electromagnetic waves backscattering from certain bodies. Práce ÚRE--ČSAV, № 17, Praha, 1961.\par
76. Chytil B. Depolarisation by randomly spaced scatterers. Práce ÚRE--ČSAV, № 20, Praha, 1961.\par
77. Chytil B. Polarisation-dependent scattering cross-sections. Práce ÚRE--ČSAV, № 21, Praha, 1961.\par
78. Beckmann P. The depolarisation of electromagnetic waves by inclined planes. Práce ÚRE--ČSAV, № 19, Praha, 1961.\par
79. Lindroth K. Reflection of electromagnetic waves from thin metal strips. Trans. Roy. Inst. of Technol., Stockholm, Sweden, № 91, 1955.\par
80. Hallén E. Theoretical investigations into the transmitting and receiving qualities of antennae. Nova Acta Roy. Soc. Sci., Upsala (4), vol. 11, Ser. IV, 1938.\par
81. Вайнштейн Л. А. Волны тока в тонком цилиндрическом проводнике. I. Ток и импеданс передающего вибратора. ЖТФ, 29, № 6, 673--688, 1959.\par
82. Вайнштейн Л. А. Волны тока в тонком цилиндрическом проводнике. II. Ток в пассивном вибраторе и излучение передающего вибратора. ЖТФ, 29, № 6, 689--699, 1959.\par
83. Вайнштейн Л. А. Волны тока в тонком цилиндрическом проводнике. III. Вариационный метод и его применение к теории идеального и импедансного проводов. ЖТФ, 31, № 1, 29--44, 1961.\par
84. Вайнштейн Л. А. Волны тока в тонком цилиндрическом проводнике. IV. Входной импеданс вибратора и точность формул. ЖТФ, 31, № 1, 45--50, 1961.\par
85. Леонтович М. А. и Левин М. Л. К теории возбуждения колебаний в вибраторах антенн. ЖТФ, 14, № 9, 481--506, 1944.\par
86. Van Vleck J. H., Bloch F. and Hamermesh M. Theory of radar reflection from wires or thin metallic strips. Journ. Appl. Phys., 18, № 3, 274--294, 1947.\par
87. Tai C. T. Electromagnetic back-scattering from cylindrical wires. Journ. Appl. Phys., 23, № 8, 903--916, 1952.\par
88. Капица П. Л., Фок В. А. и Вайнштейн Л. А. Симметричные колебания идеально проводящего полого цилиндра конечной длины. ЖТФ, 29, № 10, 1188--1205, 1959.\par
89. Говорун Н. Н. Численное решение интегрального уравнения первого рода для плотности тока в антенне --- теле вращения. «Журнал вычислительной математики и математической физики», 1, № 4, 664--679, 1961.\par
90. Hallén E. Exact Solution of the Antenna Equation. Transactions of the Royal Institute of Technology, Stockholm, Sweden, № 183, 1961.\par
91. Duncan R. H. and Hinchey F. A. Cylindrical Antenna Theory. J. Research Nat. Bur. Standards 64D, 569, 1960.\par
92. Вайнштейн Л. А. Статические задачи для полого цилиндра конечной длины (II. Численные результаты; III. Приближенные формулы). ЖТФ, 32, № 10, 1962.\par}
\end{scanpage}
